% -------------------------------------------------------
% ASSIGNED TO: \textcolor{green}{Person 4}
% -------------------------------------------------------
% INSTRUCTION: Paste existing Sections 4.4 through 4.4.3 here.
% Convert first-person to third person throughout.
% -------------------------------------------------------

% [Paste Section 4.4 SAM Regularisation here]
% [Paste Section 4.4.1 Intervention Objective here]
% [Paste Section 4.4.2 SAM Optimisation Mechanism here]
% [Paste Section 4.4.3 Expected Effect on Generalisation here]
The second stage of the experimental pipeline retrains the DeepCCDS architecture using the Sharpness-Aware Minimisation (SAM) 
optimiser to reduce the generalisation gap and improve loss landscape smoothness.

\subsubsection{Intervention Objective}
SAM is applied to the same DeepCCDS architecture, dataset splits, preprocessing pipeline, and training configuration used in the baseline 
model, ensuring that the optimisation strategy is the sole variable distinguishing the two models. This controlled substitution 
isolates the effect of loss landscape geometry on three downstream outcomes: the generalisation gap, the smoothness of the loss 
landscape as measured through visualisation, and the quality of rules extracted by the evolutionary algorithms in Sections 4.2 and 4.3.

\subsubsection{SAM Optimisation Mechanism}
SAM modifies the standard training objective by introducing a neighbourhood-based perturbation step. Rather than minimising the loss 
at the current parameters $\theta$ alone, the optimiser evaluates the worst-case loss within a local $\epsilon$-ball around the current weights, 
effectively minimising the loss at perturbed parameters $\theta + \epsilon$. This forces the model to favour solutions that remain stable under 
small parameter perturbations, penalising sharp curvature in the loss surface and encouraging convergence toward flat minima 
\cite{Foret2021,Li2018}.

\subsubsection{Expected Effect on Generalisation}
The baseline model is characterised by high training performance, lower validation performance, and a large generalisation gap, indicating 
that the model has memorised training-specific noise rather than learning stable biological patterns. By converging on flatter minima, 
the SAM-regularised model is expected to reduce the training-validation discrepancy, producing more balanced performance across both 
sets. Following \cite{Springer2024}, this improvement in generalisation is expected to produce more balanced feature representations, 
which in turn should yield extracted rules with fewer spurious conditions and higher biological coherence than those produced by the 
baseline model.
